% Contenu partagé — exercices et corrigés.
%
% Le sort des corrigés dépend uniquement de l'option solutions= du paquet.
% Ce fichier est identique dans les trois modes.

\chapter{Exercices et corrigés}%
\label{chap:exercises}
\minitoc%

\begin{quotation}
    L'option \texttt{solutions=none} ne compose aucun corrigé --- c'est la
    version étudiante. \texttt{solutions=inline} les place sous l'énoncé.
    \texttt{solutions=end} les reporte en fin de partie, avec un renvoi dans
    chaque sens. Le fichier source, lui, ne change pas.
\end{quotation}

%---------------------------------------------------------------------------
\section{Formes d'exercice}%
\label{sec:exforms}

\begin{exercise}
    Un exercice simple, sans étiquette ni barème.
\solution
    Son corrigé.
\end{exercise}

\begin{exercise}[label=ex:matrices, points=4]
    Un exercice étiqueté et barémé. On le cite par \verb|\ref{ex:matrices}|.
    Calculer $e^{tA}$ pour $A = \begin{pmatrix} 0 & 1 \\ -1 & 0\end{pmatrix}$.
\solution
    On part de la définition $e^{tA} = \sum_{n \ge 0} \frac{t^n}{n!} A^n$ et de
    l'identité $A^2 = -I$, d'où
    \[
        e^{tA} = \begin{pmatrix} \cos t & \sin t \\ -\sin t & \cos t\end{pmatrix}.
    \]
\end{exercise}

L'exercice précédent est l'exercice~\ref{ex:matrices}.

\begin{exercise}[nosolution]
    Un exercice sans corrigé, quel que soit le mode~: la clef
    \texttt{nosolution} l'exclut individuellement, et le titre ne porte alors
    aucun renvoi.
\end{exercise}

%---------------------------------------------------------------------------
\section{Questions et sous-questions}%
\label{sec:questions}

\begin{exercise}[points=6]
    \begin{question}
        Une première question, numérotée dans l'exercice.
    \end{question}
    \begin{question}
        Une deuxième question.
        \begin{subquestion}
            Une sous-question.
        \end{subquestion}
        \begin{subquestion}
            Une autre sous-question.
        \end{subquestion}
    \end{question}
    \begin{question}
        Une troisième question.
    \end{question}
\solution
    Le compteur de questions repart à zéro dans le corrigé~; \verb|\newquestion|
    permet de le faire avancer à la main quand le corrigé ne suit pas l'énoncé
    question par question.
    \begin{question}
        Corrigé de la première question.
    \end{question}
    \begin{question}
        Corrigé de la deuxième.
    \end{question}
\end{exercise}

%---------------------------------------------------------------------------
\section{Correction hors boîte}%
\label{sec:correction}

Dans un TD ou un sujet d'examen, l'énoncé n'est pas toujours encadré. Le bloc
\texttt{correction} suit alors le texte et disparaît en mode
\texttt{solutions=none}.

\begin{correction}
    Une correction brève, au fil du texte.
\end{correction}
